(1) 束缚体电流密度是 \(\mathbf J_b=\nabla\times\mathbf M=0\)，因为磁化在空间内部均匀。圆柱侧面外法向为 \(\hat r\)，所以束缚面电流 \(\mathbf K_b=\mathbf M\times\hat r=M\hat\phi\)，等价于理想无限长螺线管的周向面电流。
(2) 对跨越柱侧面的矩形安培回路，磁场的轴向跃变满足 \(B_{z,\rm in}-B_{z,\rm out}=\mu_0M\)。在无限长、无外场且远处场趋零的条件下，\(B_{\rm out}=0\)，\(B_{\rm in}=\mu_0M\)；内部 \(\mathbf H=\mathbf B/\mu_0-\mathbf M=0\)，外部同为零。
(3) 真空磁场能密度是 \(u_B=B^2/(2\mu_0)\)。内场均匀，横截面积为 \(\pi a^2\)，故单位轴向长度磁场能量 \(U_B'=\pi a^2(\mu_0M)^2/(2\mu_0)=\pi\mu_0a^2M^2/2\)。该能量为非负、尺寸量纲为焦耳每米，且 \(M=0\) 时自然归零。
(4) 柱体沿轴向平移不变且具有旋转对称性，等效束缚电流恰与无穷长螺线管相同，因而外磁场为零。若柱长有限，端面会产生退磁场及漏磁；或者磁化沿横向时面电流角分布改变，便不能直接套用本轴向无限长解。 特别地，磁场能的计算仅限于题设理想空间磁感应强度分布；若要求真实磁性材料的完整热力学自由能，还必须计入材料内部的磁化自由能，不能直接把两者等同。 对理想轴向无限圆柱，上述结论也可简写为 \(B_{mathrm{out}}=0\)，其成立取决于无限长对称性。
